\documentclass[11pt]{article}
\usepackage[a4paper,margin=25mm]{geometry}
\usepackage{amsmath,amssymb,amsthm,booktabs,hyperref,xurl}
\hypersetup{hidelinks}
\setlength{\emergencystretch}{2em}
\newtheorem{lemma}{Lemma}\newtheorem{theorem}[lemma]{Theorem}
\newcommand{\Q}{\mathbb Q}\newcommand{\Z}{\mathbb Z}
\title{A retained aggregate shift supplies the common contact spacing}
\author{\(\mathbb Z\mathbb Q\) Diophantine Project}
\date{23 November 2026}
\begin{document}
\maketitle

\begin{abstract}
The arithmetic-progression contact construction uses a free spacing
coordinate. We identify an existing aggregate shift whose two-adic
valuation is exactly the common lower bound for every contact list.
One sixteenth of this shift separates up to sixteen indexed blocks.
It depends only on fixed difference points, so it remains fixed when
the reference index is subsequently adjusted at comparison primes.
The stored spacing is an affine combination of two retained contact
coordinates in both the general and specified charts. This removes
one coordinate without changing the number or degree of the rows.
\end{abstract}

\section{The fixed difference points and the order of choices}

Use the simultaneous preparation of \cite{Shared} and the affine
contact progression of \cite{Progression}. The scalar label in this
section is denoted \(a_s\); it is unrelated to the Pell parameter
of the outer source. Write \(K=v_2(a_s)\). The original preparation
in both charts \cite{Scalar,Shared}
chooses \(64\mid h\) and a positive exact depth
\(v_2(a_s/h)\); hence \(K\ge2\) already holds. No fixed upper
bound on this depth is imposed. Its compatibility with the odd-place
annihilation orders is precisely the retained simultaneous preparation.

For each integer input \(X_i\), put \(\beta_i=51X_i^2+2504\).
The common Lucas sequence is
\[
 U_0=0,\quad U_1=1,\quad U_{n+1}=(a_s+1)U_n-U_{n-1},
 \qquad t=U_{a_s},\quad b_s=U_{a_s-1}/(a_s-1).
\]
The fixed difference points are
\[
 Q_{1,i}=[2\beta_i]R_0,\quad Q_2=[a_s]R_0,
 \quad Q_3=[t]R_0,\quad Q_4=[(b_s-1)/2]R_0.
\]
They are fixed before the contact primes are chosen and do not change
when the reference index \(h\) is later retuned. Their coordinates,
and the input-dependent scalar normalizers \(R_i\), determine all
aggregate shifts below.

\section{The common Lucas point strictly dominates at two}

Normalize the valuation on \(\Q_2(\sqrt2)\) by \(v_2(2)=1\).
The established formal calculation on the source curve \cite{Formal} gives, for
every nonzero even integer \(n\),
\begin{equation}\label{formal}
 v_2\bigl(x([n]R_0)\bigr)=-2v_2(n)-1.
\end{equation}
If \(x=x_0+x_1\sqrt2\), its two summands have valuations in
distinct cosets modulo \(\Z\). Since \eqref{formal} is integral,
the real component has exactly that valuation and the imaginary
component has valuation at least that value.

\begin{lemma}\label{Lucas}
For \(a_s\in6\Z_{>0}\) and \(K=v_2(a_s)\ge2\),
\[
 v_2(t)=2K+1,\qquad v_2(b_s-1)=K.
\]
\end{lemma}
\begin{proof}
Put \(P=a_s+1\). The recurrence gives
\(U_3=P^2-1=a_s(a_s+2)\), of valuation \(K+1\), and
\(V_3=P^3-3P\equiv2\pmod4\) for the companion trace sequence.
The addition identity gives \(U_{3m}/U_3=U_m(V_3)\).
For a recurrence parameter congruent to two modulo four, an odd
index has odd \(U\)-value, and every companion trace value is
congruent to two modulo four. The doubling identity
\(U_{2j}=U_jV_j\) therefore contributes exactly one factor of
two at each doubling. Thus \(v_2(U_{3m}/U_3)=v_2(m)\).
Taking \(m=a_s/3\) gives the first formula.

For the second, write \(a_s-1=6l-1\), and regard \(U_n\)
as an integer polynomial in \(P\). At \(P=1\), the recurrence
and its derivative give
\(U_{6l-1}(1)=-1\), \(U'_{6l-1}(1)=-2l\).
Taylor expansion at \(P=1\), with \(P-1=a_s\), yields
\(U_{a_s-1}(a_s+1)\equiv-1\pmod{2^{K+1}}\).
Thus its difference from \(a_s-1\) has valuation exactly \(K\).
Division by the odd integer \(a_s-1\) proves the result.
\end{proof}

The elementary integer-label and scalar-normalizer calculations \cite{Scalar} give
\(v_2(\beta_i)\le3\) and \(v_2(R_i)\in\{2,3\}\).
Consequently the valuations of the real components in the four root
contributions are
\begin{center}
\begin{tabular}{lc}
\toprule
Contribution & Real-component valuation\\
\midrule
\(x(Q_{1,i})\) & at least \(-9\)\\
\(R_i x(Q_2)\) & at least \(-2K+1\)\\
\(x(Q_2)\) in the specified chart & \(-2K-1\)\\
\(x(Q_3)\) & exactly \(-4K-3\)\\
\(x(Q_4)\) & exactly \(-2K+1\)\\
\bottomrule
\end{tabular}
\end{center}
For \(K\ge2\), the \(Q_3\) term has strictly smaller valuation
than every other contribution. Every imaginary component has at
least the same lower bound \(-4K-3\).

\section{The retained shift separates every contact list}

In the general chart the aggregate is
\[
 T_i=7+\frac14\operatorname{Tr}
       \bigl(x(Q_{1,i})+R_i x(Q_2)+x(Q_3)+x(Q_4)\bigr).
\]
The preceding strict dominance gives
\begin{equation}\label{Tval}
                         v_2(T_i)=-4K-4.
\end{equation}
Its two internally disjoint contact lists are translates of the fixed
contact set by \(0\) and \(T_i\). The fixed contact constants are
two-adically integral, so every contact position belongs to
\(2^{-4K-4}\Z_2\).

In the specified chart write
\[
 S_{0,i}+S_{1,i}\sqrt2=x(Q_{1,i})+x(Q_2)+x(Q_3)+x(Q_4).
\]
Its two successive shifts are
\[
 T_i=13+S_{0,i}/2,\qquad V_i=17+(S_{0,i}-S_{1,i})/2.
\]
Again \(v_2(T_i)=-4K-4\), while
\(v_2(V_i)\ge-4K-4\). All fixed contact constants have valuation
at least \(-1\). Thus the three lists, translated by \(0,T_i,T_i+V_i\),
have the same common lower bound.

\begin{lemma}[Spacing from one fixed aggregate]\label{separation}
Choose the input block labelled \(b\) as the root block and put
\(D=T_b/16\). For distinct block indices \(i,j\) in
\(\{0,\ldots,15\}\), the translated contact lists
\(\mathcal A_i-iD\) and \(\mathcal A_j-jD\) are disjoint.
\end{lemma}
\begin{proof}
The valuation of every difference of two old contact positions is
at least \(-4K-4\). But \(v_2(D)=-4K-8\), and
\(v_2(i-j)\le3\), so
\[
                  v_2((i-j)D)\le-4K-5.
\]
These quantities cannot be equal. This proves the assertion,
including repeated integer inputs.
\end{proof}

The choice of \(D\) is made entirely from the fixed difference
points. It therefore stays fixed during every later reference-index
congruence and approximation step. Apply the original finite-family
supplier to the disjoint translated lists, with the same enlargement
of finite avoidance as in \cite{Progression}. No moving coefficients
occur. The common slope \(\lambda\) is still free to satisfy the
selected-prime congruences and the retained local and real limits.
At all those places \(iD\) is a fixed rational translation, so the
old asymptotic \(q(\lambda+iD+u)=\kappa\lambda^{-2}
+O(\lambda^{-3})\) and every global-gauge argument apply unchanged.

\section{Affine substitution in the complete quadratic charts}

Write \(t=\texttt{int\_b\_2}\), \(t'=\texttt{int\_b\_68}\),
and \(d=\texttt{contact\_common\_spacing}\). In the general
chart \(t=\lambda/\rho\), \(t'=(\lambda+T_b)/\rho\),
and the spacing coordinate is \(d=D/\rho\). In the specified
chart \(t=\lambda\), \(t'=\lambda+T_b\), and \(d=D\).
Thus in both charts the exact substitution is
\begin{equation}\label{substitution}
                              d=(t'-t)/16.
\end{equation}
The respective retained aggregate rows have constant terms fourteen
and twenty-six, and exactly give these physical shifts.

\begin{theorem}
In either complete affine-contact table with block indices contained
in \(\{0,\ldots,15\}\), substitute \eqref{substitution} everywhere
and delete \(d\). The resulting complete system has one fewer
variable, the same number of quadratic residuals, and the same outer
integer-input projection under the retained arithmetic interfaces.
\end{theorem}
\begin{proof}
For soundness assign \(d\) by \eqref{substitution}; every old row is
then exactly the substituted row. No division by a variable occurs,
and every inherited gauge guard remains. Conversely choose the
simultaneous integer witness using Lemma~\ref{separation}. The original
progression construction then satisfies \eqref{substitution} and
fills every other coordinate as before. Because the substitution is
affine with rational constant coefficients, all rows remain quadratic;
clearing constant denominators gives integer polynomial coefficients.
\end{proof}

The historical candidate using a scalar secant coordinate remains
unproved because that coordinate moves with the reference. The present
argument uses an invariant of the fixed difference points and does not
assert any property of that earlier candidate.

\section{Complete tables and verification}

Apply the affine substitution to the frozen product-elimination
predecessors in
\nolinkurl{undecidability/pell_product_v99/}. The general family has
thirteen blocks indexed by \(0,\ldots,8,10,\ldots,13\); the
specified family has eleven blocks indexed by
\(0,\ldots,8,10,11\). Both lie within the stated sixteen-index
range. The complete censuses are
\begin{center}
\begin{tabular}{lrrrr}
\toprule
Interface & Old variables & New variables & Old rows & New rows\\
\midrule
General & 972 & 971 & 826 & 826\\
Specified & 1146 & 1145 & 996 & 996\\
\bottomrule
\end{tabular}
\end{center}
All source equations and integer predicates remain unchanged. The
ordinary-\(abc\) interface uses its existing fixed offsets and is
unaffected by this transformation. The general fixed parameter and
the specified arithmetic interfaces retain precisely their existing
scopes; no extra prime-supplier or source-universality hypothesis is
introduced.

The deterministic builder and the independent audit are
\begin{verbatim}
python undecidability/aggregate_spacing_v100/build_aggregate_spacing.py
python undecidability/aggregate_spacing_v100/verify_aggregate_spacing.py
\end{verbatim}
The checker imports no builder. It expands the affine substitution
independently, checks all fourteen or twenty-two changed rows and
all 812 or 974 unchanged rows, verifies constant row multipliers,
integer polynomial coefficients, full coordinate support, degree,
the retained aggregate and gauge rows, metadata, and hashes. It also
checks seventy-five exact modular Lucas fixtures and all 120 positive
differences among sixteen indices. The finite fixtures support the
arithmetic proof without replacing its quantified statements.

\begin{thebibliography}{9}
\small\setlength{\itemsep}{2pt}
\bibitem{Shared} \(\mathbb Z\mathbb Q\) Diophantine Project,
\emph{Shared quartic checkpoint}, simultaneous scalar preparation and
fixed difference-point construction.
\par\nolinkurl{undecidability/papers/shared_quartic_checkpoint.tex}.
\bibitem{Progression} \(\mathbb Z\mathbb Q\) Diophantine Project,
\emph{A common arithmetic progression of contact arguments and shared
coefficients of their centered powers}.
\par\nolinkurl{undecidability/general_contact_v94/affine_contact_progressions.tex}.
\bibitem{Formal} \(\mathbb Z\mathbb Q\) Diophantine Project,
\emph{Paired contact quartic construction}, formal-parameter calculation.
\par\nolinkurl{papers/quartic_existence/paired_contact56_20261119_v81.tex}.
\bibitem{Scalar} \(\mathbb Z\mathbb Q\) Diophantine Project,
\emph{Scalar point charts for the general and specified interfaces}.
\par\nolinkurl{papers/quartic/scalar_point163_179_20261120_v82.tex}.
\end{thebibliography}
\end{document}
